\chapter*{Abstract}
\begin{otherlanguage}{english}
Music has cultural and social dimensions that may not be equally accessible to deaf or hard-of-hearing people, particularly when musical interfaces rely predominantly on hearing. This work presents the development and technical validation of a multisensory digital musical interface for music accessibility. The architecture uses the Akai MPD218 MIDI controller as the input interface and the \textit{velocity} parameter as a digital representation of the intensity of the gesture performed on the pads. MIDI events are processed in Python, converted into PWM commands, and transmitted through UDP over Wi-Fi to an ESP32-based embedded unit. In its final configuration, the wearable unit integrates a Vibracall 1027 vibrotactile motor driven by a MOSFET switching stage and a high-brightness LED, as well as a self-contained power system composed of two lithium-ion cells in a 2S configuration, a BMS module, and a DC--DC converter. Validation included simulations of the actuation stage, communication tests, control-logic tests, and experimental characterization of both the input and the actuator. In the UDP test, all 100 transmitted packets were recorded by the ESP32, with no losses in the evaluated sequence. In the 39~g body-drop experiment, 60 drops were performed from heights between 0.5~cm and 5~cm; the association between height and MIDI \textit{velocity} yielded a Spearman coefficient of 0.944 ($p<0.001$). In the acoustic characterization of the motor, six PWM levels from 80 to 255 were evaluated in three repetitions, with mean frequencies ranging from 91.5~Hz to 227.9~Hz and a Spearman coefficient of 0.988 ($p<0.001$). With the embedded power supply, voltages between approximately 3.26~V and 3.29~V were measured at the ESP32 supply rail with Wi-Fi active. Integrated operation confirmed the functional chain between MIDI interaction, computer processing, wireless communication, and vibrotactile and visual responses. The results demonstrate the technical feasibility of the architecture as an experimental platform, without establishing, at this stage, perceptual effectiveness for deaf or hard-of-hearing users.

\noindent\textbf{Keywords:} music accessibility; MIDI; vibrotactile feedback; multisensory interaction; ESP32.

\end{otherlanguage}

\clearpage
